\documentclass[11pt,a4paper]{ltjsarticle}

\usepackage[margin=24mm]{geometry}
\usepackage{amsmath,amssymb,bm}
\usepackage{booktabs,tabularx,array}
\usepackage{xcolor}
\usepackage{hyperref}

\hypersetup{
  colorlinks=true,
  linkcolor=blue!50!black,
  urlcolor=blue!60!black,
  pdftitle={5-step Latent ResfusionからStable Diffusionへの直接handover}
}
\renewcommand{\arraystretch}{1.18}

\newcommand{\E}{\mathbb{E}}
\newcommand{\N}{\mathcal{N}}

\title{5-step Latent Resfusionから\\
Stable Diffusionへの直接handover}
\author{}
\date{\today}

\begin{document}
\maketitle

\section{目的}

構成は次のとおりとする．
\[
\begin{aligned}
x
&\longrightarrow \text{VAE encoder}
\longrightarrow z_0
\longrightarrow \text{ADJSCC encoder} \\
&\longrightarrow \text{wireless channel}
\longrightarrow \text{ADJSCC decoder}
\longrightarrow \hat z_0
\longrightarrow \text{5-step Latent Resfusion} \\
&\longrightarrow \text{Stable Diffusion}
\longrightarrow \text{VAE decoder}.
\end{aligned}
\]

Resfusionは原論文\cite{resfusion}と公式実装\cite{resfusion-code}どおりの
residual noiseだけを学習する．追加のresidual head，bridge network，
score loss，統計loss，画像lossは使わない．

推論時には，Resfusionを0回から5回まで実行した各状態を
Stable Diffusionへ直接渡せるようにする．ここで「0回」は
Resfusionの初期状態をそのまま渡す場合，「5回」は5回すべてdenoiseした後を表す．

\section{学習用latent pair}

画像 \(x\) からStable Diffusionのscaled VAE latentを作る．
\[
  z_0=s\,E_{\mathrm{VAE}}(x),
  \qquad s=0.18215.
\]
このlatentをADJSCCで伝送する．
\[
\begin{aligned}
  u &= E_{\mathrm{JSCC}}(z_0),\\
  y &= \operatorname{Channel}(u),\\
  \hat z_0 &= D_{\mathrm{JSCC}}(y).
\end{aligned}
\]
したがって，Resfusionのpaired dataは
\[
  (\text{degraded latent},\text{clean latent})
  =
  (\hat z_0,z_0)
\]
であり，prior residualは
\begin{equation}
  R=\hat z_0-z_0
  \label{eq:residual}
\end{equation}
となる．VAEと学習済みADJSCCは固定し，Latent Resfusionだけを学習する．

scaled latentはRGB画像ではないため，公式RGB実装にある
\([0,1]\to[-1,1]\) 変換と \([-1,1]\) clampは行わない．

\section{論文どおりのResfusion学習}

\subsection{forward process}

\[
  \alpha_t=1-\beta_t,
  \qquad
  \bar\alpha_t=\prod_{i=0}^{t}\alpha_i.
\]
以下では，質問中の \(\alpha_t\) をcumulative alpha
\(\bar\alpha_t\) として表記する．

Resfusionのforward stateを
\begin{equation}
  v_t
  =
  \sqrt{\bar\alpha_t}\,z_0
  +
  \left(1-\sqrt{\bar\alpha_t}\right)R
  +
  \sqrt{1-\bar\alpha_t}\,\epsilon,
  \qquad
  \epsilon\sim\N(0,I)
  \label{eq:forward}
\end{equation}
とする．式\eqref{eq:residual}を代入すると
\begin{equation}
  v_t
  =
  \left(2\sqrt{\bar\alpha_t}-1\right)z_0
  +
  \left(1-\sqrt{\bar\alpha_t}\right)\hat z_0
  +
  \sqrt{1-\bar\alpha_t}\,\epsilon
  \label{eq:smooth}
\end{equation}
となる．

\subsection{学習target}

Resfusion U-Netが予測するresidual noiseは
\begin{equation}
  \eta_t
  =
  \epsilon
  +
  \frac{
    \left(1-\sqrt{\alpha_t}\right)
    \sqrt{1-\bar\alpha_t}
  }{\beta_t}R
  \label{eq:resnoise}
\end{equation}
である．

各batchで
\[
  t\sim\operatorname{Uniform}\{0,1,2,3,4\},
  \qquad
  \epsilon\sim\N(0,I)
\]
をsampleし，式\eqref{eq:forward}で \(v_t\) を作る．
networkへの入力は
\[
  (v_t,\hat z_0,t)
\]
だけであり，4チャネルの \(v_t\) と4チャネルの \(\hat z_0\) を結合するため，
U-Netの入力は8チャネル，出力は4チャネルとする．

\subsection{loss}

学習lossは原論文のMSEだけとする．
\begin{equation}
\boxed{
  \mathcal L_{\mathrm{Resfusion}}
  =
  \E_{z_0,\hat z_0,t,\epsilon}
  \left[
    \left\|
      \eta_\theta(v_t,\hat z_0,t)-\eta_t
    \right\|_2^2
  \right]
}
\label{eq:loss}
\end{equation}

使用しないものは次のとおりである．
\[
  \mathcal L_R,\quad
  \mathcal L_{\mathrm{bridge}},\quad
  \mathcal L_{\mathrm{score}},\quad
  \mathcal L_{\mathrm{stat}},\quad
  \mathcal L_{\mathrm{image}}.
\]
また，SNR専用embeddingや追加headも設けない．channel状態の違いは
degraded condition \(\hat z_0\) 自体からnetworkへ伝わる．

\section{5-step Resfusion}

\subsection{truncated schedule}

公式復元設定と同じ
\[
  T=12,\qquad \text{LinearProScheduler}
\]
を使用する．truncated strategyにより
\[
  \sqrt{\bar\alpha_{t_\star}}=\frac12
\]
となる点でscheduleが切られ，実際に使用する時刻は
\[
  t_{\mathrm R}=4,3,2,1,0
\]
の5個になる．したがってU-Net評価は5回である．

Resfusion初期状態は
\begin{equation}
\begin{aligned}
  v_4
  &=
  \sqrt{\bar\alpha^{\mathrm R}_4}\,\hat z_0
  +
  \sqrt{1-\bar\alpha^{\mathrm R}_4}\,\epsilon\\
  &=
  0.5\,\hat z_0+\sqrt{0.75}\,\epsilon
\end{aligned}
\label{eq:init}
\end{equation}
とする．

各reverse stepは
\begin{equation}
  v_{t-1}
  =
  \frac{1}{\sqrt{\alpha_t}}
  \left(
    v_t
    -
    \frac{\beta_t}{\sqrt{1-\bar\alpha_t}}
    \eta_\theta(v_t,\hat z_0,t)
  \right)
  +
  \sqrt{\tilde\beta_t}\,\xi,
  \qquad \xi\sim\N(0,I)
  \label{eq:reverse}
\end{equation}
である．最後のstepでは \(\xi=0\) とする．

\subsection{0--5 stepの定義}

Resfusionを完了した回数を \(k\) とする．
\[
\begin{array}{ccl}
  k=0 &: & v_4 \quad\text{（初期化直後）},\\
  k=1 &: & v_3 \quad\text{（\(t_{\mathrm R}=4\) を1回denoise後）},\\
  k=2 &: & v_2,\\
  k=3 &: & v_1,\\
  k=4 &: & v_0,\\
  k=5 &: & v_{-1}\quad\text{（5回denoise後のclean推定）}.
\end{array}
\]
この6状態のいずれからでもStable Diffusionへhandoverする．

\section{Stable Diffusion timestepとの対応}

\subsection{対応規則}

ResfusionとStable Diffusionではnoise scheduleが異なるため，
Resfusionのstep番号をそのままStable Diffusionへ渡してはいけない．

質問中の
\[
  \left(1-\sqrt{\bar\alpha_t}\right)\hat z_0
  +
  \sqrt{1-\bar\alpha_t}\,\epsilon
\]
に現れる係数をnoise levelの識別に使う．Resfusion状態 \(k\) と
Stable Diffusion raw timestep \(\tau\) の距離を
\begin{equation}
\begin{aligned}
d(k,\tau)
={}&
\left[
  \left(1-\sqrt{\bar\alpha^{\mathrm R}_k}\right)
  -
  \left(1-\sqrt{\bar\alpha^{\mathrm{SD}}_\tau}\right)
\right]^2\\
&+
\left[
  \sqrt{1-\bar\alpha^{\mathrm R}_k}
  -
  \sqrt{1-\bar\alpha^{\mathrm{SD}}_\tau}
\right]^2
\end{aligned}
\label{eq:distance}
\end{equation}
とし，
\begin{equation}
\boxed{
  \tau_k
  =
  \arg\min_{\tau\in\{0,\ldots,999\}}d(k,\tau)
}
\label{eq:mapping}
\end{equation}
でStable Diffusion側の時刻を選ぶ．

これは，signal/noise比
\[
  \rho_t
  =
  \frac{\sqrt{1-\bar\alpha_t}}
       {\sqrt{\bar\alpha_t}}
\]
が最も近い時刻を選ぶことと実質的に同じである．

\subsection{本プロジェクトでの対応表}

Stable Diffusion v1の
\[
  T_{\mathrm{SD}}=1000,\qquad
  \texttt{linear\_start}=0.00085,\qquad
  \texttt{linear\_end}=0.012
\]
を用いる．
\path{ldm/modules/diffusionmodules/util.py} と同じく，
端点の平方根を線形補間してから二乗した \(\beta_\tau^{\mathrm{SD}}\) を使う．

\begin{table}[htbp]
  \centering
  \small
  \caption{Resfusionの0--5 denoise回数とStable Diffusion raw timestep}
  \label{tab:mapping}
  \begin{tabular}{crrrrr}
    \toprule
    完了回数 \(k\) &
    Resfusion state &
    \(\bar\alpha^{\mathrm R}\) &
    \(\sqrt{\bar\alpha^{\mathrm R}}\) &
    \(\sqrt{1-\bar\alpha^{\mathrm R}}\) &
    SD raw \(\tau_k\) \\
    \midrule
    0 & \(v_4\)    & 0.250000 & 0.500000 & 0.866025 & 520 \\
    1 & \(v_3\)    & 0.310431 & 0.557163 & 0.830403 & 475 \\
    2 & \(v_2\)    & 0.575518 & 0.758629 & 0.651523 & 309 \\
    3 & \(v_1\)    & 0.833902 & 0.913182 & 0.407552 & 146 \\
    4 & \(v_0\)    & 0.991667 & 0.995825 & 0.091287 & 9 \\
    5 & \(v_{-1}\) & 1.000000 & 1.000000 & 0.000000 & 0 \\
    \bottomrule
  \end{tabular}
\end{table}

よってhandover先は
\[
\boxed{
  k=0,1,2,3,4,5
  \quad\longrightarrow\quad
  \tau_k=520,475,309,146,9,0
}
\]
となる．

\(k=5\) はnoise 0であるが，Stable Diffusionの既存scheduleには
厳密なnoise 0のraw timestepがない．そのため最も近い \(\tau=0\) を使う．

\section{直接handover推論}

handoverでは補正，re-noise，residual subtractionを一切行わない．
Resfusionの現在状態をそのままStable Diffusion latentとして使う．
\begin{equation}
\boxed{
  z^{\mathrm{SD}}_{\tau_k}
  \leftarrow
  v^{(k)}
}
\label{eq:direct}
\end{equation}
ここで \(v^{(k)}\) はResfusionを \(k\) 回denoiseした状態である．

推論手順は次のとおりである．
\begin{enumerate}
  \item ADJSCC decoderから \(\hat z_0\) を得る．
  \item 式\eqref{eq:init}で \(v_4\) を作る．
  \item 選択した \(k\) 回だけ式\eqref{eq:reverse}を実行する．
  \item 表\ref{tab:mapping}から \(\tau_k\) を得る．
  \item 式\eqref{eq:direct}によりcurrent stateを直接Stable Diffusionへ渡す．
  \item Stable Diffusionをraw timestep \(\tau_k\) から0まで実行する．
  \item 得られたscaled latentをVAE decoderで画像へ戻す．
\end{enumerate}

\subsection{疑似コード}

\begin{verbatim}
z_hat0 = adjscc_decode(received_signal)

v = 0.5 * z_hat0 + sqrt(0.75) * randn_like(z_hat0)

sd_timestep = [520, 475, 309, 146, 9, 0]

for k in range(6):
    if k == selected_handover_step:
        z_out = stable_diffusion_reverse(
            latent=v,
            start_raw_timestep=sd_timestep[k]
        )
        image_out = vae_decode_scaled(z_out)
        break

    # k=0,...,4 のときだけ次のResfusion stepを実行
    t_resfusion = 4 - k
    eta_hat = resfusion_unet(v, z_hat0, t_resfusion)
    v = resfusion_reverse_step(v, eta_hat, t_resfusion)
\end{verbatim}

\section{最小構成のまとめ}

\begin{itemize}
  \item Resfusionは公式設定と同じ5回のdenoiseを行う．
  \item 学習targetはresidual noise \(\eta_t\) だけである．
  \item lossは式\eqref{eq:loss}のMSEだけである．
  \item U-Netは8チャネル入力，4チャネル出力である．
  \item 0--5回の各状態を直接Stable Diffusionへ渡せる．
  \item 対応raw timestepは \(520,475,309,146,9,0\) である．
  \item handover時にtensorの補正や追加noiseは行わない．
\end{itemize}

\begin{thebibliography}{9}

\bibitem{resfusion}
Z. Shi et al.,
\newblock ``Resfusion: Denoising Diffusion Probabilistic Models for
Image Restoration Based on Prior Residual Noise,''
\newblock \emph{NeurIPS 2024}.
\newblock \url{https://arxiv.org/abs/2311.14900}

\bibitem{resfusion-code}
NKICSL,
\newblock ``Official implementation of Resfusion.''
\newblock \url{https://github.com/nkicsl/Resfusion}

\end{thebibliography}

\end{document}
